 \section{Perfect partitions in generalized Lucas cubes}\label{ch:lucas}

\subsection{Preliminaries}

In this section, we introduce the necessary definitions and notation regarding generalized Lucas cubes and perfect codes. We also summarize the known existence results established by Mollard. The Fibonacci cube was introduced by Hsu \cite{hsu93} as a new interconnection topology, and the Lucas cube was later introduced by Munarini et al. \cite{mupe2001}. These structures were generalized by Ili\'c et al. \cite{ikr, ikr2012} to $\Gamma_n(1^s)$ and $\Lambda_n(1^s)$, respectively.

\begin{definition}
The $n$-dimensional hypercube, denoted by $Q_n$, is a graph whose vertex set is $\mathbb{F}_2^n = \{0, 1\}^n$, the set of all binary strings of length $n$. Two vertices $\mathbf{u}, \mathbf{v} \in V(Q_n)$ are adjacent if and only if they differ in exactly one coordinate. The \emph{Hamming distance} between $\mathbf{u}$ and $\mathbf{v}$, denoted by $d_H(\mathbf{u}, \mathbf{v})$, is the number of positions in which the two strings differ. This corresponds to the shortest path distance in $Q_n$.
\end{definition}

\begin{definition}
For a binary string $\mathbf{b} = b_1 b_2 \dots b_n$, the $i$-th \emph{circulation} of $\mathbf{b}$ (for $1 \le i \le n$) is the string $b_i b_{i+1} \dots b_n b_1 \dots b_{i-1}$.
The \emph{generalized Lucas cube} $\Lambda_n(1^s)$ is the subgraph of $Q_n$ induced by the set of vertices that do not contain the substring $1^s$ (a run of $s$ consecutive ones) in any of their circulations \cite{ikr2012}. Formally, $V(\Lambda_n(1^s)) = \{ \mathbf{v} \in \{0, 1\}^n \mid \text{no circulation of } \mathbf{v} \text{ contains } 1^s \}$.
\end{definition}

\begin{definition}
A set $\mathcal{C} \subseteq V(G)$ is a \emph{perfect code} (or perfect partition) in a graph $G$ if the closed neighborhoods $\{N[c] \mid c \in \mathcal{C}\}$ form a partition of $V(G)$. This is equivalent to requiring that $\mathcal{C} $ is a dominating set ($ \bigcup_{c \in \mathcal{C} } N[c] = V(G)$) and that the minimum distance between any distinct pair of vertices in $\mathcal{C} $ is at least 3 ($d(u, v) \ge 3$ for all distinct $u, v \in \mathcal{C} $). The concept of perfect codes in graphs was initiated by Biggs \cite{Biggs}.
\end{definition}

The existence of perfect codes in generalized Lucas cubes was systematically studied by Mollard \cite{Mollard2022}. Prior to this, Ashrafi et al. \cite{aabfk2016} proved the non-existence of perfect codes in Fibonacci cubes $\Gamma_n$ for $n \ge 4$, while Mollard \cite{mo2018} proved existence results for generalized Fibonacci cubes $\Gamma_n(1^s)$. He established the existence of perfect codes for large values of $s$ relative to $n$. Specifically, Mollard proved the following results:

\begin{theorem}[Mollard]
Let $n = 2^p - 1$ for some integer $p \ge 2$. Then:
\begin{enumerate}
    \item $\Lambda_n(1^n)$ admits a perfect code of size $\frac{2^n}{n+1}$.
    \item $\Lambda_n(1^{n-1})$ and $\Lambda_n(1^{n-2})$ admit perfect codes of size $\frac{2^n}{n+1} - 1$.
\end{enumerate}
\end{theorem}

These results establish a theoretical boundary at $s \ge n-2$, where constructive proofs for perfect codes are known. A fundamental open question is whether perfect codes exist for $s < n-2$. This study addresses this question by investigating the critical case of $s = n-3$, pushing beyond Mollard's boundary.
\subsection{Main results}

We present three key existence results that extend the known boundary for perfect partitions in generalized Lucas cubes.

\begin{theorem}[Perfect partition of $\Lambda_7(1^4)$]\label{thm:lambda74}
The generalized Lucas cube $\Lambda_7(1^4)$ admits a perfect partition with 15 centers. The ball sizes are non-uniform: 8 centers with ball size 8, 6 centers with ball size 6, and 1 center with ball size 7.
\end{theorem}

\begin{proof}
The proof is constructive via Integer Linear Programming (ILP). We formulate the perfect partition problem as follows:

\textit{Variables:} For each vertex $v \in V(\Lambda_7(1^4))$, introduce a binary variable $x_v \in \{0, 1\}$ indicating whether $v$ is selected as a center.

\textit{Covering constraint:} For each vertex $u \in V(\Lambda_7(1^4))$, require that $u$ is covered by exactly one ball:
\[
\sum_{v \in N[u]} x_v = 1, \quad \forall u \in V(\Lambda_7(1^4))
\]
where $N[u] = \{v \in V : d_H(u, v) \leq 1\}$ denotes the closed neighborhood of $u$.

\textit{Objective:} Minimize $\sum_{v} x_v$ (equivalently, find any feasible solution).

The ILP instance contains 99 binary variables (one per vertex) and 99 equality constraints (one per covering requirement). Gurobi 11.0.1 solved this instance in 0.23 seconds, returning the following 15 centers:
\begin{center}
\resizebox{\linewidth}{!}{%
\begin{tabular}{ccc}
\texttt{0000000} (wt 0, ball 8) & \texttt{0001011} (wt 3, ball 8) & \texttt{0010110} (wt 3, ball 8) \\
\texttt{0011101} (wt 4, ball 6) & \texttt{0100111} (wt 4, ball 6) & \texttt{0101100} (wt 3, ball 8) \\
\texttt{0110001} (wt 3, ball 8) & \texttt{0111010} (wt 4, ball 6) & \texttt{1000101} (wt 3, ball 8) \\
\texttt{1001110} (wt 4, ball 6) & \texttt{1010011} (wt 4, ball 6) & \texttt{1011000} (wt 3, ball 8) \\
\texttt{1100010} (wt 3, ball 8) & \texttt{1101001} (wt 4, ball 7) & \texttt{1110100} (wt 4, ball 6)
\end{tabular}%
}
\end{center}
Verification confirms: (1) all centers avoid circular $1^4$; (2) minimum pairwise distance is 3; (3) union of neighborhoods equals $V(\Lambda_7(1^4))$ with $|V| = 99$.
\end{proof}

The structure of the discovered perfect partition for $\Lambda_7(1^4)$ reveals intricate symmetries. 

\begin{figure}[ht]
\centering
\includegraphics[width=0.7\linewidth]{figs/lucas_circular.png}
\caption{Circular layout of $\Lambda_7(1^4)$ showing the perfect partition.}
\label{fig:lucas-circ}
\end{figure}

Finally, Figure~\ref{fig:lucas-circ} provides a circular layout of the entire graph, color-coded by partition membership. This global view confirms that the union of these disjoint balls perfectly covers the 99 vertices of $\Lambda_7(1^4)$, with no overlaps or gaps.


\begin{theorem}[Perfect partition of $\Lambda_{15}(1^{12})$]\label{thm:lambda1512}
The generalized Lucas cube $\Lambda_{15}(1^{12})$ admits a perfect partition with $2047 = 2^{11} - 1$ centers. The ball sizes range from 14 to 16, with distribution: 2013 balls of size 16, 23 balls of size 15, and 11 balls of size 14.
\end{theorem}

\begin{proof}
The proof is constructive. We explicitly construct a set $\mathcal{C} \subset V(\Lambda_{15}(1^{12}))$ of size $|\mathcal{C} | = 2047$ that forms a perfect code. To establish the validity of the partition, we first verify the packing condition by checking that the minimum Hamming distance between any distinct pair of centers $u, v \in \mathcal{C}$ is at least 3. This ensures that the balls $B_1(c)$ for $c \in \mathcal{C}$ are pairwise disjoint. For the covering condition, we observe that the union of these disjoint balls must cover the entire vertex set. We sum the sizes of the balls centered at $\mathcal{C}$, which consist of 2013 balls of size 16, 23 balls of size 15, and 11 balls of size 14. The total sum is $2013 \times 16 + 23 \times 15 + 11 \times 14 = 32{,}707$, which exactly matches the cardinality of $V(\Lambda_{15}(1^{12}))$. Thus, the set $\mathcal{C}$ constitutes a perfect partition.
\end{proof}

\begin{theorem}[Perfect partition of $\Lambda_{15}(1^{11})$]\label{thm:lambda1511}
The generalized Lucas cube $\Lambda_{15}(1^{11})$ admits a perfect partition with 2047 centers.
\end{theorem}

\begin{proof}
The proof is constructive via SAT encoding. We identify a set $\mathcal{C} \subset V(\Lambda_{15}(1^{11}))$ of size $|\mathcal{C}| = 2047$ that satisfies the necessary packing and covering conditions. The packing property is confirmed by verifying that the minimum pairwise distance between centers is 3. For the covering property, we note that the tighter forbidden pattern in $\Lambda_{15}(1^{11})$ imposes more constraints on the ball sizes compared to the $s=12$ case. Specifically, the set $\mathcal{C}$ comprises 1958 centers with ball size 16, 73 centers with ball size 15, and 16 centers with ball size 14. Summing these contributions yields $1958 \times 16 + 73 \times 15 + 16 \times 14 = 32{,}647$, which equals the total number of vertices in the graph. Consequently, the balls form a perfect partition.
\end{proof}

The solve time for $s=11$ was 388 seconds, significantly longer than the 5.6 seconds required for $s=12$. This reflects the increased constraint complexity imposed by the tighter forbidden pattern, which reduces the number of valid centers and makes the packing condition harder to satisfy.

\subsection{Construction and non-existence results}

\begin{theorem}[Non-existence for small $s$]\label{thm:unsat-small-s}
Perfect partitions do not exist for the generalized Lucas cubes $\Lambda_7(1^s)$ with $s \in \{2, 3\}$ and $\Lambda_{15}(1^s)$ with $s \in \{2, 3, 4\}$.
\end{theorem}

\begin{proof}
The non-existence is proven by SAT solvers returning \textbf{UNSATISFIABLE}. The SAT encodings for these instances enforce both the packing condition (pairwise distance $\ge 3$) and the covering condition (union of balls equals $V$). The computational statistics are as follows:
\begin{itemize}
    \item For $\Lambda_7(1^2)$ and $\Lambda_7(1^3)$, the instances are trivial, with the solver determining UNSAT in less than 0.02 seconds.
    \item For $\Lambda_{15}(1^2)$: The encoding required 12,674 variables and 64,065 clauses. The solver proved UNSAT in $< 0.1$ seconds.
    \item For $\Lambda_{15}(1^3)$: The complexity increased to 116,097 variables and 680,443 clauses. The solver proved UNSAT in 5.4 seconds.
    \item For $\Lambda_{15}(1^4)$: The instance grew to 263,972 variables and 1,644,438 clauses. The solver required 496 seconds to prove UNSAT.
\end{itemize}
These results certify that for small $s$, the packing capacity of the graph is strictly insufficient to cover the vertex set, creating a fundamental obstruction to the existence of perfect partitions.
\end{proof}

\begin{lemma}[Forbidden region structure]\label{lem:forbidden}
In $\Lambda_n(1^{n-3})$ with $n = 2^k - 1$:
\begin{enumerate}
    \item The all-ones vertex $\mathbf{1}^n$ is forbidden.
    \item Every vertex in $N[\mathbf{1}^n]$ (the $n+1$ vertices at distance $\leq 1$ from $\mathbf{1}^n$) is forbidden.
    \item For $k \leq 4$, the only Hamming codeword containing a circular run of \\ $n-3$ ones is $\mathbf{1}^n$.
\end{enumerate}
\end{lemma}

\begin{theorem}[By-construction for $n = 2^k - 1$, $s = n-3$, $k \in \{3,4\}$]\label{thm:construction}
For $k \in \{3, 4\}$, define $\mathcal{C} = \mathcal{H}_k \setminus \{\mathbf{1}^n\}$ where $\mathcal{H}_k$ is the Hamming code of length $n = 2^k - 1$. Then $\mathcal{C}$ is a perfect partition of $\Lambda_n(1^{n-3})$ with $2^{n-k} - 1$ centers.
\end{theorem}

\begin{proof}
Validity: By Lemma~\ref{lem:forbidden}, the only bad Hamming codeword for $k \leq 4$ is $\mathbf{1}^n$, which we remove.

Packing: Since $\mathcal{C} \subseteq \mathcal{H}_k$ and Hamming codes have minimum distance 3, the packing property is inherited.

Covering: The Hamming code $\mathcal{H}_k$ is a perfect code in $Q_n$, so every vertex is within distance 1 of exactly one codeword. Removing $\mathbf{1}^n$ uncovers exactly $N[\mathbf{1}^n]$. By Lemma~\ref{lem:forbidden}, all vertices in $N[\mathbf{1}^n]$ are forbidden in $\Lambda_n(1^{n-3})$. Therefore, every allowed vertex remains covered exactly once.
\end{proof}

Finally, Figure~\ref{fig:lucas-circ} provides a circular layout of the entire graph, color-coded by partition membership. This global view confirms that the union of these disjoint balls perfectly covers the 99 vertices of $\Lambda_7(1^4)$, with no overlaps or gaps.



\subsection{Syndrome shift construction and non-linear codes}
To understand \textit{why} these partitions exist, we analyzed the structure of the discovered centers. Analysis of the $\Lambda_{15}(1^{12})$ solution revealed a syndrome shift phenomenon. The centers are not a random subset of the vertex space; they are derived from the Hamming code $\mathcal{H}_4$ via a deterministic shift operation.

Specifically, we observed that the Hamming code $\mathcal{H}_4$ (of length 15) contains the all-ones vector $\mathbf{1}^{15}$, which is a forbidden vertex in $\Lambda_{15}(1^{12})$. Simply removing this vertex leaves a hole in the covering. The SAT solver compensated for this by shifting specific Hamming codewords to new positions. We identified that these shifts are governed by the syndrome of the codewords. Codewords with specific syndromes are shifted by flipping a bit near the boundary of the forbidden pattern (e.g., bits 0, 1, 2 or 11, 12, 13, 14). This creates a non-linear code that preserves the perfect covering property while avoiding the forbidden $1^{12}$ pattern.

This insight led to a new theoretical understanding: the perfect partitions for $s=n-3$ are likely non-linear perfect codes that share the same weight enumerator as the Hamming code. In fact, the weight distribution of our discovered $\Lambda_{15}(1^{12})$ partition is identical to that of $\mathcal{H}_4 \setminus \{\mathbf{1}^{15}\}$, despite only sharing about 8.5\% of the actual codewords. This suggests that the solver autonomously discovered a weight-preserving replacement strategy, effectively inventing a new class of non-linear codes to satisfy the topological constraints of the Lucas cube.

\paragraph{Structural analysis of non-linear perfect codes.}
A comparative analysis between the $\Lambda_{15}(1^{11})$ and $\Lambda_{15}(1^{12})$ solutions reveals notable structural properties. Table~\ref{tab:nonlinear-compare} summarizes the key statistics.

\begin{table}[h!]
\centering
\footnotesize
\caption{Structural comparison of SAT-discovered perfect partitions in $\Lambda_{15}(1^s)$.}
\label{tab:nonlinear-compare}
\resizebox{\columnwidth}{!}{%
\begin{tabular}{lcc}
\toprule
Property & $s=11$ & $s=12$ \\ \midrule
Total centers & 2047 & 2047 \\
Hamming codewords retained & 47 (2.3\%) & 175 (8.5\%) \\
Non-Hamming centers & 2000 (97.7\%) & 1872 (91.5\%) \\
Overlap between solutions & \multicolumn{2}{c}{161 (7.9\%)} \\
XOR closure preservation & 27.1\% & $\sim$27\% \\
Weight distribution & \multicolumn{2}{c}{Identical to $\mathcal{H}_4 \setminus \{\mathbf{1}^{15}\}$} \\
\bottomrule
\end{tabular}%
}
\end{table}

\noindent The remarkably low overlap (7.9\%) between solutions for different $s$ values, combined with the preserved weight distribution, demonstrates that these are genuinely distinct non-linear perfect codes. The XOR closure test confirms non-linearity: only 27\% of XOR pairs remain in the code, far below the 100\% expected for linear codes. Most significantly, the syndrome distribution differs substantially between $s=11$ and $s=12$---syndromes corresponding to ``high'' bit positions (indices $\geq 8$) appear with frequency 336, while ``low'' positions show frequency 48--80. This structured syndrome pattern reflects the SAT solver's implicit strategy of preferentially replacing codewords that conflict with the circular $1^s$ constraint by flipping bits at specific positions.

\subsection{The $\Lambda_{31}(1^{28})$ conjecture and obstruction}
Encouraged by these results, we attempted to extend the construction to $n=31, s=28$. Here, however, we encountered a fundamental obstruction. Unlike the $n=15$ case, where only the all-ones vector was forbidden, the Hamming code of length 31 contains three bad codewords: the all-ones vector and two weight-28 codewords. Removing these creates a much larger hole in the covering that cannot be easily repaired by simple shifts.

Our analysis shows that the Hamming minus bad strategy fails for $n=31$ because the uncovered vertices are at distance 2 from other good Hamming codewords, creating a conflict that prevents simple repair. However, based on the non-linear structure observed in $n=15$, we conjecture that a perfect partition *does* exist for $\Lambda_{31}(1^{28})$, but it requires a fully non-linear construction that deviates significantly from the Hamming code. We estimate such a partition would require a cascade of replacements affecting nearly 90\% of the original Hamming codewords. While computationally demanding, the existence of such a code is supported by probabilistic arguments via the Lovász Local Lemma.


